\apendice{Anexo de sostenibilización curricular}

\section{Introducción}

La sostenibilidad suele asociarse principalmente a aspectos medioambientales, aunque en realidad también incluye dimensiones sociales, económicas y tecnológicas. Durante el desarrollo de este Trabajo Fin de Grado he podido comprobar que muchas de las decisiones adoptadas en un proyecto tecnológico pueden contribuir de forma directa o indirecta a una utilización más eficiente y responsable de los recursos disponibles.

El proyecto desarrollado consiste en un sistema SIEM ligero orientado a la monitorización de eventos de seguridad dentro de un entorno simulado de alta seguridad. Aunque a primera vista pueda parecer un proyecto centrado exclusivamente en aspectos técnicos, durante su desarrollo se han aplicado diversos principios relacionados con la sostenibilidad.

Uno de los aspectos más relevantes ha sido la utilización de tecnologías de código abierto. Herramientas como Python, SQLite, FastAPI, Streamlit, Git o GitHub permiten desarrollar soluciones funcionales sin necesidad de adquirir licencias de software comerciales. Esto favorece la accesibilidad tecnológica y facilita que pequeñas organizaciones o instituciones educativas puedan disponer de herramientas avanzadas sin realizar grandes inversiones económicas.

Otro aspecto importante es la eficiencia en el uso de los recursos hardware. Desde el inicio del proyecto se planteó el desarrollo de un SIEM ligero capaz de ejecutarse en equipos con recursos limitados, incluyendo plataformas como Raspberry Pi. Este enfoque permite reducir el consumo energético y prolongar la vida útil de los equipos utilizados, contribuyendo a un uso más sostenible de la tecnología.

La sostenibilidad también está relacionada con la seguridad y la continuidad de los servicios. Los sistemas de monitorización permiten detectar incidencias, fallos o situaciones anómalas antes de que se conviertan en problemas de mayor gravedad. En este sentido, el proyecto contribuye a mejorar la supervisión de infraestructuras críticas mediante mecanismos de correlación de eventos y detección de anomalías que facilitan la identificación temprana de posibles riesgos.

Desde el punto de vista social, el desarrollo del proyecto ha permitido adquirir competencias relacionadas con el trabajo autónomo, la resolución de problemas y el aprendizaje continuo. A lo largo del desarrollo ha sido necesario analizar diferentes alternativas tecnológicas, tomar decisiones de diseño y adaptarse a nuevas herramientas y metodologías. Estas capacidades resultan fundamentales para el ejercicio profesional responsable dentro del ámbito de la ingeniería informática.

Asimismo, la utilización de herramientas de inteligencia artificial como apoyo al desarrollo ha permitido mejorar la productividad y agilizar determinadas tareas técnicas y documentales. No obstante, durante el proyecto se ha mantenido una actitud crítica respecto a la información generada por estas herramientas, revisando y validando siempre los resultados obtenidos antes de incorporarlos al sistema.

Por último, considero que este trabajo me ha permitido comprender mejor la importancia de diseñar soluciones tecnológicas eficientes, mantenibles y escalables. La sostenibilidad no debe entenderse únicamente como una cuestión medioambiental, sino también como la capacidad de desarrollar sistemas que aprovechen adecuadamente los recursos disponibles, puedan mantenerse a largo plazo y aporten valor real a las organizaciones y a la sociedad.

Los aspectos abordados durante el desarrollo del proyecto pueden relacionarse especialmente con el Objetivo de Desarrollo Sostenible 9 (Industria, innovación e infraestructura), al promover el desarrollo de soluciones tecnológicas innovadoras basadas en software libre y arquitecturas eficientes. Asimismo, el proyecto guarda relación con el Objetivo de Desarrollo Sostenible 16 (Paz, justicia e instituciones sólidas), ya que contribuye al ámbito de la seguridad y la monitorización de infraestructuras sensibles mediante herramientas que facilitan la detección temprana de incidentes y la mejora de los mecanismos de supervisión.

En conclusión, el desarrollo de este Trabajo Fin de Grado ha contribuido a reforzar mi comprensión de los principios de sostenibilidad aplicados al ámbito tecnológico, mostrando cómo decisiones relacionadas con la arquitectura del sistema, las herramientas utilizadas y la gestión de los recursos pueden tener un impacto positivo desde el punto de vista económico, social y medioambiental.